\chapter{A Bicycle Ride with Uncle Ken}
Kissing the girl while sharing the bicycle with her is no easy task, but I managed it when I was thirteen and my cousin Melanie was fourteen. Of course, we both fell off in the process, and landed in one of Granny's flower beds where we were well cushioned by her nasturtiums.
\par
I was a clumsy boy, always falling off bicycles, and Cousin Melanie was teaching me to ride properly, making me sit on the front seat while she guided the infernal machine from the carrier seat. The kiss was purely expermental. I had not kissed a girl before, and as Cousin Melanie seemed eminently kissable, I thought I'd start with her. I waited until we were stationary, and she was instructing me on the intricacies of the cycle chain, and then I gave her a hurried kiss on the cheek.
\par
She was so startled that she fell backwards, talking me and the bicycle with her.
\par
Later, she reported me to Granny, who said, 'We'll have to keep an eye on that boy. He's showing signs fo a dissolute nature.'
\par
'What's dissolute, Uncle Ken?' I asked my favourite uncle.
'It means you're going to the dogs. You're not supposed to kiss your cousin.'
\par
'Can I kiss other girls?'
\par
'Only if they are willing.'
\par
'Did you ever kiss a girl, Uncle Ken?'
\par
Uncle Ken blushed. 'Er . . . well . . . a long time ago.'
\par
'Tell me about it.'
\par
'Another time.'
\par
'Now tell me now. How old were you?'
\par
'About twenty.'
\par
'And how old was she?'
\par
'A bit younger.'
\par
'And what happened?'
\par
'We went cycling together. I was living in Agra, when your grandfather worked there with the railways. Daisy's father was an engine driver. But she didn't like engines-they left her covered in soot. Everyone had a bicycle in those days, only the very rich had cars. And the cars could not keep up with the bicycles.
\par
'We lived in the cantonment, where the roads were straight and wide. Daisy and I went on cycle rides to Fatehpur Sikri and Secunderabad and, of course, the Taj Mahal, and one evening we saw the Taj Mahal by moonlight and it made us very romantic, and when I saw her home, we kissed under the Ashoka trees.'
'I didn't know you were so romantic, Uncle Ken. Why didn't you marry Daisy?'
\par
'I didn't have a job. She said she'd wait until I got one, but after two years, she got tired of waiting. She married a ticket inspector.'
\par
'Such a sad story', I said. 'And you still don't have a job.'
\par
Uncle Ken had been through various jobs-private tutor, salesman, shop assistant, hotel manager (until he brought about the closure of the hotel) and cricket coach, this last on the strength of bearing a close resemblance to Bruce Hallam-but at present he was unemployed, and only too ready to put his vast experience of life at my disposal.
\par
Not only did he teach me to ride a bicycle but also accompanied me on cycle rides around Dehra and along the lanes and country roads outside the town.
\par
A bicycle provides its rider with a great amount of freedom. A car will take you further, but the fact that you're sitting in a confined space detracts from the freedom of the open spaces and unfamiliar roads. On a cycle, you can feel the breeze on your face, small the mango trees in blossom, slow down and gaze at the buffaloes wading in their ponds, or just stop anywhere and get down and enjoy a cup of tea or a glass of sugar cane juice. Footslogging takes time and cars are too fast-everything whizzes past you before you can take a second look-and car drivers hate having to stop; they are intent only on reaching their destinations in good time. But a bicycle is just right for someone who likes to take a leisurely look at the world as well as to give the world a chance to look at him.
\par
Uncle Ken and I had some exhilarating bicycle rides during my winter holidays, and the most memorable of these was our unplanned visit to a certain 'Rest Home' situated on the outskirts of the town. It isn't there now, so don't go looking for it.
\par
We had cycled quite far that day, and were tired and thirsty. There was no sign of a tea shop on that particular road, but when we arrived at the open gate of an impressive building with a signboard saying 'Rest and Recuperation Centre', we presumed it was a hotel or hostelry fo sorts and rode straight into the premises. There was an extensive lawn to one side, surrounded by neat hedges and flowering shrubs. A number of people were strolling about on the lawn; some were sitting on benches; one or two were straddling a wall, talking to themselves; another was standing alone, singing to a non-existent audience. Some were Europeans; a few were Indians.
\par
We left our cycles on the porch and went in search of refreshment. A lady in a white sari gave us cool water from a surahi and told us we could wait on a bench just outside their office. But Uncle Ken said we'd prefer to meet some of the guests and led me across the lawn to where the singer was practising his notes. He was a florid gentleman, heavily built.
\par
'Do you like my singing?' he asked, as we came up.
\par
'Wonderful!' exclaimed Uncle Ken. 'You sing like Caruso.'
\par
'I am Caruso!' affirmed the tenor and let rip the opening notes of a famous operatic aria. 'Your tiny hand is frozen,' he sang, although it was an unusually warm day.
\par
We hurried on a met an elegant couple who were parading up and down the lawn, waving their hands to an invisible crowd.
\par
'Good day to you, gentlemen,' said a flamboyant individual. 'You're the ambassadors from Sweden, I suppose.'
\par
'If you so wish,' said Uncle Ken gallantly. 'And I have the honour of speaking to-?'
\par
'The Emperor Napoleon, of course.'
\par
'Of course. And this must be the Empress Joséphine.' Uncle Ken bowed to the lady beside him.
\par
'Actually, it's Marie Walewska,' said Napoleon. 'Joséphine is indisposed today.'
\par
I was beginning to feel like Alice at the Mad Hatter's tea party, and began tugging at Uncle Ken's coat sleeve, whispering that we were getting late for lunch.
\par
A turbaned warrior with a tremendous moustache loomed in front of us. 'I'm Prithviraj Chauhan,' he announced. 'And I invite you to dinner at my palace.'
\par
'Everyone's royalty here,' observed Uncle Ken. 'It's such a privilege to be with you.'
\par
'Me too,' I put in nervously.
\par
'Come with me, boy, and I'll introduce you to the others.' Prithviraj Chauhan took me by the hand and began guiding me across the lawn. 'There are many famous men and women here. That's Marco Polo over there. He's just back from China. And if you don't care for Caruso's singing, there's Tansen, under that tamarind tree. Tamarind leaves are good for the voice, you know that, of course. And that fashinable gentleman there, he's Lord Curzon, who used to be a viceroy. He's talking to the Sultan of Marrakesh. Come along, I'll introduce you to them . . . You're the young prince fo Denmark, aren't you?'
\par
'Hamlet himself,' said Uncle Ken with a wink.
\par
Before I could refute any claims to royalty, we were intercepted by a white-coated gentleman accompanied by a white-coated assistant. They looked as through they were in charge.
\par¨
'And what are you doing here, young man?' asked the senior of the two.
\par
'I'm with my uncle,' I said, gesturing towards Uncle Ken, who approached and gave the incharge an affable handshake.
\par
'And you must be Dr Freud,' said Uncle Ken. 'I must say this is a jolly of a sort of place.'
\par
'Actually, I'm Dr Goel. You must be the new patient we were expecting. But they should have sent you over with someone  a little older than this boy. Never mind, come along to the office and we'll have you admitted.'
\par
Uncle Ken and I both protested that we were not potential patients but had entered the grounds by mistake. We had our bicycles to prove it! But Dr Goel was having nothing of this deception. He and his assistant linked arms with Uncle Ken and marched him off to the office, while I trailed behind, wondering if I should get on my bicycle and rush back to the Granny with the terrible news that Uncle Ken had been incarcerated in a lunatic asylum.
\par
Just then, an ambulance arrived with the real patient, a school principal suffering from a persecution complex. He kept shouting that he was perfectly sane, and that his entire staff had plotted to have him put away. This might well have been true, as the staff was there in force to make sure he did not escape.
\par
Dr Goel apologized to Uncle Ken. Uncle Ken apologized to Dr Goel. The good doctor even accompanied us to the gate. He shook hands with Uncle Ken and said, 'I have a feeling we'll see you here again.' He looked hard at my uncle and added, 'I think I've seen you before, sir. What did you say your name was?'
\par
'Bruce Hallam,' said Uncle Ken mischievously, and rode away before they charged their minds and kept him in their 'Rest Home'.
